% =============================================================================
% Figure 1E-3: Module 1E typical application schematic
% =============================================================================
% Pico 2 W -> AD9742 (12 data bits + clock) -> 25-ohm termination -> 5th-order
% Butterworth reconstruction filter -> 1k input resistor -> AD8056 op-amp ->
% SP3T impedance switch (50/600/10k) -> BNC.
%
% Design representation note: the AD9742 has two complementary current outputs
% (IOUTA, IOUTB) and the AD8056 is wired as a differential-to-single-ended
% converter in the actual KiCad schematic. For figure readability, this
% high-level functional view shows the signal flow as single-ended (one DAC
% output line, one termination, one signal through the filter). The full
% differential topology is documented in the design doc section 3 and lives
% in the KiCad schematic.
%
% Pin label convention: show GP0..GP2, an ellipsis, and GP11 to indicate the
% 12-bit bus rather than enumerating every pin. CLK shown separately, with
% vertical clearance from the IC title at the bottom of the box.
%
% Compile:  pdflatex 1e_typical_app.tex
% Convert:  pdftocairo -svg 1e_typical_app.pdf
% =============================================================================
\documentclass[border=8pt]{standalone}
%
% =============================================================================
% pmvb-figures.sty -- Poor Man's Validation Bench figure style sheet
% =============================================================================
% Defines the house style for all TikZ figures in PMVB design documents.
%
% Usage in figure source:
%     \documentclass[border=4pt]{standalone}
%     \usepackage{pmvb-figures}
%
% Two color modes:
%     \pmvbDarkMode   -- dark background (default; matches SDD HTML theme)
%     \pmvbLightMode  -- white background (for printed portfolio PDFs)
% =============================================================================
\NeedsTeXFormat{LaTeX2e}
\ProvidesPackage{pmvb-figures}[2026/05/06 PMVB figure styles]
%
% --- Required packages ------------------------------------------------------
\RequirePackage{tikz}
% NOTE: dropped the siunitx option. Our figures use \pmvblabel{} for unit
% annotations rather than circuitikz's SI value formatting, and siunitx.sty
% isn't always available on bare TeX Live installs.
\RequirePackage[american]{circuitikz}
\RequirePackage{xcolor}
\RequirePackage{etoolbox}
%
% --- TikZ libraries ---------------------------------------------------------
\usetikzlibrary{arrows.meta}
\usetikzlibrary{positioning}
\usetikzlibrary{shapes.geometric}
\usetikzlibrary{shapes.misc}
\usetikzlibrary{fit}
\usetikzlibrary{backgrounds}
\usetikzlibrary{calc}
\usetikzlibrary{decorations.pathreplacing}
\usetikzlibrary{patterns}
%
% =============================================================================
% Color palette (FMCW dark theme, matches SDD HTML)
% =============================================================================
\definecolor{pmvbBg}{HTML}{0d1117}
\definecolor{pmvbFg}{HTML}{e6edf3}
\definecolor{pmvbMuted}{HTML}{94a3b8}
\definecolor{pmvbBlue}{HTML}{3b82f6}
\definecolor{pmvbBlueFill}{HTML}{1f2937}
\definecolor{pmvbGreen}{HTML}{10b981}
\definecolor{pmvbGreenFill}{HTML}{1e293b}
\definecolor{pmvbAmber}{HTML}{f59e0b}
\definecolor{pmvbRed}{HTML}{ef4444}
\definecolor{pmvbViolet}{HTML}{a78bfa}
%
\definecolor{pmvbBgLight}{HTML}{ffffff}
\definecolor{pmvbFgLight}{HTML}{0f172a}
\definecolor{pmvbMutedLight}{HTML}{475569}
\definecolor{pmvbBlueFillLight}{HTML}{dbeafe}
\definecolor{pmvbGreenFillLight}{HTML}{d1fae5}
%
% =============================================================================
% Mode selectors
% =============================================================================
\newcommand{\pmvbDarkMode}{%
  \pagecolor{pmvbBg}%
  \color{pmvbFg}%
}
%
\newcommand{\pmvbLightMode}{%
  \colorlet{pmvbBg}{pmvbBgLight}%
  \colorlet{pmvbFg}{pmvbFgLight}%
  \colorlet{pmvbMuted}{pmvbMutedLight}%
  \colorlet{pmvbBlueFill}{pmvbBlueFillLight}%
  \colorlet{pmvbGreenFill}{pmvbGreenFillLight}%
  \pagecolor{pmvbBgLight}%
  \color{pmvbFgLight}%
}
%
\AtBeginDocument{\pmvbDarkMode}
%
% =============================================================================
% Node styles
% -----------------------------------------------------------------------------
% IMPORTANT: each \tikzset{...} block below must be a single argument with NO
% blank lines inside, otherwise \par is inserted and pgfkeys errors out.
% =============================================================================
%
% --- Subsystem (system-level block: a board, a chip, a whole module) ---
\tikzset{
  pmvb subsystem/.style={
    rectangle,
    rounded corners=2pt,
    draw=pmvbBlue,
    line width=0.8pt,
    fill=pmvbBlueFill,
    text=pmvbFg,
    minimum width=28mm,
    minimum height=12mm,
    align=center,
    font=\sffamily\small,
    inner sep=4pt,
  },
}
%
% --- IC / hardware part (functional-level block) ---
\tikzset{
  pmvb ic/.style={
    rectangle,
    rounded corners=1pt,
    draw=pmvbGreen,
    line width=0.6pt,
    fill=pmvbGreenFill,
    text=pmvbFg,
    minimum width=24mm,
    minimum height=10mm,
    align=center,
    font=\sffamily\footnotesize,
    inner sep=3pt,
  },
}
%
% --- Internal block (inside an IC's functional block diagram) ---
\tikzset{
  pmvb internal/.style={
    rectangle,
    draw=pmvbFg,
    line width=0.5pt,
    fill=pmvbBg,
    text=pmvbFg,
    minimum width=18mm,
    minimum height=7mm,
    align=center,
    font=\sffamily\scriptsize,
    inner sep=2pt,
  },
}
%
% --- Pin pad (the box-with-X at IC die boundary in datasheet diagrams) ---
\tikzset{
  pmvb pin/.style={
    rectangle,
    draw=pmvbFg,
    line width=0.4pt,
    fill=pmvbBg,
    minimum size=2.5mm,
    inner sep=0pt,
    path picture={
      \draw[pmvbFg, line width=0.3pt]
        (path picture bounding box.south west) -- (path picture bounding box.north east)
        (path picture bounding box.north west) -- (path picture bounding box.south east);
    },
  },
}
%
% --- Connector (front panel, BNC, banana jack) ---
\tikzset{
  pmvb connector/.style={
    circle,
    draw=pmvbFg,
    line width=0.6pt,
    fill=pmvbBg,
    text=pmvbFg,
    minimum size=8mm,
    font=\sffamily\scriptsize,
    inner sep=1pt,
  },
}
%
% --- Annotation (no border, small muted italic text) ---
\tikzset{
  pmvb annotation/.style={
    text=pmvbMuted,
    font=\sffamily\scriptsize\itshape,
    align=center,
  },
}
%
% --- Subsystem grouping box (used with \node[fit=...]) ---
\tikzset{
  pmvb group/.style={
    rectangle,
    rounded corners=4pt,
    draw=pmvbMuted,
    line width=0.4pt,
    dashed,
    text=pmvbMuted,
    inner sep=8pt,
    font=\sffamily\scriptsize,
  },
}
%
% =============================================================================
% Edge styles
% -----------------------------------------------------------------------------
% Edge styles encode the *kind* of signal flowing on the line.
% =============================================================================
%
\tikzset{
  pmvb digital/.style={
    -Stealth,
    draw=pmvbBlue,
    line width=0.8pt,
    font=\sffamily\scriptsize,
  },
}
%
\tikzset{
  pmvb analog/.style={
    -Stealth,
    draw=pmvbAmber,
    line width=0.8pt,
    font=\sffamily\scriptsize,
  },
}
%
\tikzset{
  pmvb power/.style={
    -Stealth,
    draw=pmvbRed,
    line width=1.0pt,
    font=\sffamily\scriptsize,
  },
}
%
\tikzset{
  pmvb control/.style={
    -Stealth,
    draw=pmvbViolet,
    line width=0.7pt,
    dashed,
    font=\sffamily\scriptsize,
  },
}
%
\tikzset{
  pmvb bus/.style={
    -Stealth,
    draw=pmvbBlue,
    line width=1.4pt,
    double=pmvbBg,
    double distance=0.6pt,
    font=\sffamily\scriptsize,
  },
}
%
% Style for the "xN" bus-width annotation (overlaid on a bus line).
% Dark fill masks the underlying double-line so the white text reads cleanly.
\tikzset{
  pmvb bus width/.style={
    fill=pmvbBg,
    rounded corners=1pt,
    inner sep=2pt,
    font=\sffamily\bfseries\scriptsize,
    text=pmvbFg,
  },
}
%
\tikzset{
  pmvb optional/.style={
    -Stealth,
    draw=pmvbMuted,
    line width=0.6pt,
    dashed,
    font=\sffamily\scriptsize,
  },
}
%
% --- Bidirectional variants ---
\tikzset{
  pmvb digital bi/.style={
    pmvb digital,
    Stealth-Stealth,
  },
}
%
\tikzset{
  pmvb analog bi/.style={
    pmvb analog,
    Stealth-Stealth,
  },
}
%
% =============================================================================
% Diagram-level convenience
% =============================================================================
\tikzset{
  pmvb figure/.style={
    >=Stealth,
    every node/.append style={font=\sffamily},
    node distance=10mm and 14mm,
  },
}
%
% =============================================================================
% Helper macros
% =============================================================================
\newcommand{\pmvblabel}[1]{\textcolor{pmvbMuted}{\sffamily\scriptsize\itshape #1}}
%
\newcommand{\pmvbsubsystem}[3]{%
  node[pmvb subsystem] (#1) {{\bfseries #2}\\\scriptsize #3}%
}
%
\newcommand{\pmvbic}[3]{%
  node[pmvb ic] (#1) {{\bfseries #2}\\\scriptsize #3}%
}
%
% NOTE on bus-width annotations: TikZ's path parser doesn't reliably expand
% \newcommand macros at mid-path positions, so for the "xN" overlay use the
% inline node-in-path syntax directly:
%     \draw[pmvb bus] (A) -- node[pmvb bus width]{$\times 16$} (B);
% The pmvb bus width style is defined above.
%
% -----------------------------------------------------------------------------
% \opamptri{name}{coord}{dir}
%   Draws an op-amp triangle. Apex points right (dir=1) or left (dir=-1).
%   Defines coordinates: <name>_p, <name>_n, <name>_out for wiring.
%   Place any label as a separate \node call after the macro.
%
%   The {coord} argument MUST include its own parentheses, e.g.
%       \opamptri{bufa}{(0.6, -0.5)}{1}
%       \opamptri{tl072}{($(m1e.center) + (4mm, -14mm)$)}{1}
%   This lets callers use either simple coords or calc expressions.
% -----------------------------------------------------------------------------
\newcommand{\opamptri}[3]{%
  \begin{scope}[shift={#2}, xscale=#3]
    \draw[pmvbFg, line width=0.6pt, fill=pmvbBg]
      (0, 0) -- (
%
% Shrink default circuitikz bipole length so resistors fit cleanly inline.
\ctikzset{bipoles/length=8mm}
%
\begin{document}
\begin{circuitikz}[pmvb figure]
  %
  % =========================================================================
  % LEFT: Pico 2 W with abbreviated GPIO label list on right edge
  % =========================================================================
  \node[pmvb ic, minimum width=24mm, minimum height=50mm, anchor=west]
    (pico) at (0, 0) {};
  %
  % Title at the BOTTOM of the box. yshift=4mm of south places it at
  % y = -2.5 + 0.4 = -2.1
  \node[font=\sffamily\bfseries\small, text=pmvbFg]
    at ([yshift=4mm]pico.south) {Pico 2 W};
  %
  % Abbreviated GPIO label list. Pin pitch 5mm; CLK row pulled higher
  % (y=-1.4) so it's clear of the title (y=-2.1).
  \coordinate (pico_gp0)  at ($(pico.east) + (0,  1.6)$);
  \coordinate (pico_gp1)  at ($(pico.east) + (0,  1.1)$);
  \coordinate (pico_gp2)  at ($(pico.east) + (0,  0.6)$);
  \coordinate (pico_dots) at ($(pico.east) + (0,  0.1)$);
  \coordinate (pico_gp11) at ($(pico.east) + (0, -0.4)$);
  \coordinate (pico_gp12) at ($(pico.east) + (0, -1.4)$);
  %
  \node[font=\sffamily\scriptsize, text=pmvbFg, anchor=east]
    at ($(pico_gp0)  + (-2mm, 0)$) {GP0};
  \node[font=\sffamily\scriptsize, text=pmvbFg, anchor=east]
    at ($(pico_gp1)  + (-2mm, 0)$) {GP1};
  \node[font=\sffamily\scriptsize, text=pmvbFg, anchor=east]
    at ($(pico_gp2)  + (-2mm, 0)$) {GP2};
  \node[font=\sffamily\scriptsize, text=pmvbFg, anchor=east]
    at ($(pico_dots) + (-2mm, 0)$) {\dots};
  \node[font=\sffamily\scriptsize, text=pmvbFg, anchor=east]
    at ($(pico_gp11) + (-2mm, 0)$) {GP11};
  \node[font=\sffamily\scriptsize, text=pmvbFg, anchor=east]
    at ($(pico_gp12) + (-2mm, 0)$) {GP12};
  %
  % =========================================================================
  % MIDDLE-LEFT: AD9742 with abbreviated DB pin label list on left edge
  % =========================================================================
  \node[pmvb ic, minimum width=28mm, minimum height=50mm]
    at (5, 0) (dac) {};
  %
  \node[font=\sffamily\bfseries\small, text=pmvbFg]
    at ([yshift=4mm]dac.south) {AD9742};
  %
  \coordinate (dac_db0)  at ($(dac.west) + (0,  1.6)$);
  \coordinate (dac_db1)  at ($(dac.west) + (0,  1.1)$);
  \coordinate (dac_db2)  at ($(dac.west) + (0,  0.6)$);
  \coordinate (dac_dots) at ($(dac.west) + (0,  0.1)$);
  \coordinate (dac_db11) at ($(dac.west) + (0, -0.4)$);
  \coordinate (dac_clk)  at ($(dac.west) + (0, -1.4)$);
  %
  \node[font=\sffamily\scriptsize, text=pmvbFg, anchor=west]
    at ($(dac_db0)  + (2mm, 0)$) {DB0};
  \node[font=\sffamily\scriptsize, text=pmvbFg, anchor=west]
    at ($(dac_db1)  + (2mm, 0)$) {DB1};
  \node[font=\sffamily\scriptsize, text=pmvbFg, anchor=west]
    at ($(dac_db2)  + (2mm, 0)$) {DB2};
  \node[font=\sffamily\scriptsize, text=pmvbFg, anchor=west]
    at ($(dac_dots) + (2mm, 0)$) {\dots};
  \node[font=\sffamily\scriptsize, text=pmvbFg, anchor=west]
    at ($(dac_db11) + (2mm, 0)$) {DB11};
  \node[font=\sffamily\scriptsize, text=pmvbFg, anchor=west]
    at ($(dac_clk)  + (2mm, 0)$) {CLK};
  %
  % VDD / VSS labels and connections
  \coordinate (dac_vdd)  at (dac.north);
  \node[font=\sffamily\scriptsize, text=pmvbFg, anchor=north]
    at ($(dac_vdd) + (0, -1mm)$) {$V_{DD}$};
  %
  \coordinate (dac_vss)  at (dac.south);
  \draw (dac_vss) to[short, *-] ++(0, -0.4) node[ground]{};
  %
  % IOUT (single-ended representation) on right edge.
  % y=0.183 matches the op-amp's .+ input (offset above center by ~1.83mm at
  % scale=0.85 with yscale=-1). With IOUT, LPF, and the 1k resistor all at
  % y=0.183, the entire input side of the analog chain is horizontal. The
  % op-amp's output sits at y=0 (center of triangle), so the output wire to
  % SP3T is also horizontal at y=0. The ~1.83mm vertical step happens inside
  % the op-amp triangle.
  \coordinate (dac_iout) at ($(dac.east) + (0, 0.183)$);
  \node[font=\sffamily\scriptsize, text=pmvbFg, anchor=east]
    at ($(dac_iout) + (-2mm, 0)$) {IOUT};
  %
  % =========================================================================
  % Pico -> AD9742: 12-bit bus (no extra text label, x12 carries the meaning)
  % + separate CLK line below
  % =========================================================================
  % Bus center: at midpoint of GP0..GP11 cluster (≈ y=0.6)
  \draw[pmvb bus] ($(pico.east) + (0, 0.6)$) --
    node[pmvb bus width]{$\times 12$}
    ($(dac.west) + (0, 0.6)$);
  %
  % CLK: separate single digital line
  \draw[pmvb digital] (pico_gp12) --
    node[above=0.5mm]{\pmvblabel{CLK}}
    (dac_clk);
  %
  % =========================================================================
  % BYPASS NETWORK above AD9742
  % =========================================================================
  % Coords raised 1cm vs initial draft: the AD9742 box (50mm tall) extends to
  % y=2.5, so the cap+ground network needs to land well above that. Junction
  % line at y=4.3, caps drop 0.7cm to ground glyphs at y=3.6 (11mm clearance
  % above the box top).
  \node[font=\sffamily\scriptsize, text=pmvbRed]
    at (5, 5.0) (vdd_label) {+3.3 V};
  %
  \coordinate (c1_top)  at (3.8, 4.3);
  \coordinate (vdd_jct) at (5.0, 4.3);
  \coordinate (c2_top)  at (6.2, 4.3);
  %
  \draw[pmvbFg, line width=0.5pt]
    (vdd_label.south) -- (vdd_jct) -- (dac_vdd);
  \draw[pmvbFg, line width=0.5pt] (c1_top) -- (c2_top);
  \fill[pmvbFg] (vdd_jct) circle (0.5mm);
  %
  \draw (c1_top) to[C, l_={\pmvblabel{C1 10\,$\mu$F}}, *-]
    ++(0, -0.7) node[ground]{};
  \draw (c2_top) to[C, l^={\pmvblabel{C2 0.1\,$\mu$F}}, *-]
    ++(0, -0.7) node[ground]{};
  %
  % =========================================================================
  % MIDDLE: 25-ohm termination + LPF (single-ended representation)
  % =========================================================================
  % Single horizontal arrow from IOUT all the way to the LPF input. The 25Ω
  % termination resistor branches off the wire at a tap point and goes down
  % to ground. A junction dot at the tap shows the connection.
  \node[pmvb ic, minimum width=22mm, minimum height=14mm,
        draw=pmvbAmber, fill=pmvbBg]
    at (10, 0.183) (lpf) {};
  \node[font=\sffamily\bfseries\scriptsize, text=pmvbFg]
    at ([yshift=-3mm]lpf.north) {5th-order};
  \node[font=\sffamily\scriptsize, text=pmvbFg]
    at (lpf.center) {Butterworth LPF};
  \node[font=\sffamily\scriptsize\itshape, text=pmvbMuted]
    at ([yshift=2mm]lpf.south) {12 MHz};
  %
  \coordinate (lpf_in)  at (lpf.west);
  \coordinate (lpf_out) at (lpf.east);
  %
  % Tap point on the IOUT-to-LPF wire (where the 25Ω branches off)
  \coordinate (term_tap) at ($(dac_iout) + (0.7, 0)$);
  %
  % Single horizontal arrow from IOUT to LPF input (one continuous line,
  % arrow head at the LPF end)
  \draw[pmvb analog] (dac_iout) -- (lpf_in);
  %
  % 25Ω resistor branches off at the tap, vertical to ground
  \draw (term_tap) to[R=\pmvblabel{25\,$\Omega$}]
    ($(term_tap) + (0, -0.7)$) node[ground]{};
  %
  % Junction dot at the tap (indicates connection between signal wire and
  % termination resistor)
  \fill[pmvbFg] (term_tap) circle (0.5mm);
  %
  % =========================================================================
  % MIDDLE-RIGHT: 1k input resistor -> AD8056 op-amp
  % =========================================================================
  % Op-amp default orientation (no yscale flip): + input at bottom-left,
  % - input at top-left. The 1k resistor must enter the + input, not the -.
  % We use yscale=-1 here to put + on TOP for visual convention; this also
  % flips - to bottom, so the feedback loop drops cleanly down from .- and
  % wraps under the triangle.
  \node[op amp, scale=0.85, yscale=-1] (opa) at (13, 0) {};
  % Title placed 6mm above op-amp center (above the triangle top vertex,
  % which sits at center+4.25mm with scale=0.85 yscale=-1). Don't use
  % opa.up — that anchor is the V+ supply pin which flips to the bottom
  % under yscale=-1.
  \node[font=\sffamily\bfseries\scriptsize, text=pmvbFg]
    at ($(opa.center) + (0, 6mm)$) {AD8056};
  %
  % 1k resistor between LPF output and op-amp + input. With yscale=-1,
  % .+ is at top-left of the triangle (above center). LPF output is at y=0.
  % Route the wire LPF.east → small jog up → 1k resistor (still horizontal-ish)
  % → op-amp .+. Acceptable to have a small vertical step here since the +
  % input is offset from y=0 by ~1.7mm at scale=0.85.
  \draw (lpf_out) to[R=\pmvblabel{1\,k$\Omega$}] (opa.+);
  %
  % Feedback wire from op-amp - input (now at bottom-left due to yscale=-1)
  % to a T-junction point past the triangle apex on the output wire.
  % Wire path: opa.- → DOWN 5mm → RIGHT past apex → UP to fb_tap.
  % This routes well clear of the triangle outline.
  \coordinate (fb_tap) at ($(opa.out) + (5mm, 0)$);
  \draw (opa.-) -- ++(0, -5mm) -| (fb_tap);
  \fill[pmvbFg] (fb_tap) circle (0.5mm);
  %
  % =========================================================================
  % RIGHT: SP3T impedance switch + BNC output
  % =========================================================================
  % SP3T center y=0 matches the rest of the analog signal chain so the
  % output wire from op-amp.out (also at y=0) is perfectly horizontal.
  \node[pmvb ic, minimum width=22mm, minimum height=42mm]
    at (16, 0) (sp3t) {};
  \node[font=\sffamily\bfseries\small, text=pmvbFg]
    at ([yshift=4mm]sp3t.south) {SP3T};
  \node[font=\sffamily\scriptsize\itshape, text=pmvbMuted]
    at ([yshift=8mm]sp3t.south) {(3 reeds)};
  %
  \node[font=\sffamily\scriptsize, text=pmvbFg, anchor=west]
    at ($(sp3t.west) + (2mm,  0.7)$) {50\,$\Omega$};
  \node[font=\sffamily\scriptsize, text=pmvbFg, anchor=west]
    at ($(sp3t.west) + (2mm,  0.0)$) {600\,$\Omega$};
  \node[font=\sffamily\scriptsize, text=pmvbFg, anchor=west]
    at ($(sp3t.west) + (2mm, -0.7)$) {10\,k$\Omega$};
  %
  % Op-amp output wire passes through the feedback tap (junction dot
  % already drawn above) before continuing to the SP3T input.
  \draw[pmvb analog] (opa.out) --
    node[above=0.5mm]{\pmvblabel{$\pm$10\,V}}
    (sp3t.west);
  %
  \node[pmvb connector, right=8mm of sp3t.east] (bnc) {BNC};
  \draw[pmvb analog] (sp3t.east) -- (bnc.west);
  %
  % =========================================================================
  % Bottom annotation
  % =========================================================================
  \node[pmvb annotation, below=10mm of dac.south]
    {C1 = 10\,$\mu$F tantalum (bulk decoupling) \textbullet{} C2 = 0.1\,$\mu$F ceramic (HF decoupling)\\
     LPF: 5th-order Butterworth, $\sim$12 MHz cutoff (L--C--L--C--L ladder, 50\,$\Omega$ characteristic)\\
     SP3T: three Coto 9007-05-01 reed relays, mutually exclusive Pico GPIO drive\\
     AD9742 actually has differential outputs (IOUTA, IOUTB); shown single-ended here for figure clarity};
  %
\end{circuitikz}
\end{document}
